% Noncircular bridges for the literal dependency-annotated Hilbert-FOL sheet.
\begin{annexitem}{annex:p2.hfa.kernels}{thm:p02-hfa-propositional-kernels}{Lifting the existing propositional basis, not assuming FOL completeness}
\begin{proof}
The A1--A3 formulas and MP on the two sheets are literally the same
schemas. Replace every propositional variable in a given finite
propositional sequence by its assigned first-order fragment formula,
commuting with \(\neg,\to\). An axiom instance remains an instance
of the identical schema, and both MP premise strings still match its
conclusion string. Give a replaced premise its singleton support, an
axiom empty support, and an MP conclusion the union. Induction over
the sequence proves the asserted lifted derivation and its support.
There is no quantified inference in the operation, even if a substituted
formula itself contains quantifiers.

For a \(\neg,\to\) tautology, the map in
Definition~\ref{def:p02-hilbert-translation} acts identically on that
fragment. The cited Theorem~\ref{thm:p02-calculi-derivability-equivalence}
with empty premises supplies an A1--A3/MP derivation. Lifting it gives
the required empty-support kernel. This is the sole imported
propositional-completeness dependency, with status \texttt{cited=Kleene1952};
it is not presented as a newly reproduced source proof. If an explicit
kernel sequence is wanted, enumerate finite A1--A3/MP sequences over
the finite propositional variables in its template until that tautology
is reached. The cited basis theorem ensures termination of this
particular proof search. Any auxiliary propositional variables in an
initial basis proof may first be replaced by one variable of the
nonempty template, leaving its endpoint unchanged. Substitution then yields the first-order
kernel. Thus using a propositional kernel is a specified finite
construction, not a new first-order inference rule.
\end{proof}
\end{annexitem}

\begin{annexitem}{annex:p2.hfa.soundness}{thm:p02-hfa-soundness}{Each literal rule and its dependency-sensitive semantic condition}
\begin{proof}
Induct over the annotated sequence, proving \(D\entails A\) for
each line \((A,D)\). Premise lines have their singleton support.
A1 and A2 are classical truth-function tautologies. For A3, if its
conclusion \(\psi\) were false, \(\neg\psi\) would be true.
The two assumed implications would make both \(\neg\varphi\)
and \(\varphi\) true, which is impossible. Thus its truth function
is also valid. These three schemas have empty support. MP uses the
two induction statements at the same structure and assignment, taking
the union of their supports.

The universal instantiation schema is valid by the universal truth
clause and the actual substitution lemma, Theorem~\ref{thm:fol-substitution}.
The existential-introduction schema is valid by the reverse calculation:
the term's value supplies a witness. Bound-variable renamings introduced
by total substitution preserve satisfaction, so no additional partial
substitutability assumption is needed.

For the universal distribution schema, suppose \(\forall x(B\to C)\)
and \(B\) hold under an assignment. For any domain element \(a\),
the updated assignment still satisfies \(B\), because
\(x\notin\FV(B)\), by formula coincidence. Its \(B\to C\)
instance is true by the universal premise, and hence \(C\) is true
under that update. This holds for every \(a\), proving \(\forall xC\).
Thus the complete schema is valid and has empty support.

Suppose Gen takes \((B,D)\) to \((\forall xB,D)\).
Fix a structure and assignment satisfying \(D\). The rule's condition
puts \(x\) outside the free variables of every member of \(D\),
so every update of \(x\) still satisfies \(D\). The induction
hypothesis gives \(B\) at each update; therefore \(\forall xB\)
holds at the initial assignment. Formulas elsewhere in the permitted
premise set but outside \(D\) are not needed in this reasoning.

Finally suppose existential elimination takes \((B\to C,D)\)
to \((\exists xB\to C,D)\). Fix a model and assignment satisfying
\(D\) and \(\exists xB\). Choose a witness \(a\) for that
existential. Updating \(x\) to \(a\) preserves every member of
\(D\) by the premise restriction. The induction hypothesis makes
\(B\to C\) true at the update, and the witness makes \(B\)
true, so \(C\) is true there. The separate condition
\(x\notin\FV(C)\) transfers it to the original assignment.
This proves the rule and exhausts the literal sheet. The proof uses
neither equality axioms nor an interpretation of equality as identity.
Every final support is contained in \(\Gamma\), proving the
arbitrary-premise soundness conclusion.
\end{proof}
\end{annexitem}

\begin{annexitem}{annex:p2.hfa.discharge}{thm:p02-hfa-discharge-interface}{Tracked discharge through Gen and existential elimination}
\begin{proof}
The A1/A2 identity derivation of \(B\to B\) is the explicit one
in Annex~\ref{annex:p2.calc.2}; all its dependencies are empty.
We strengthen the usual line induction: from a line \((C,E)\) obtain
\(B\to C\) with support contained in \(E\setminus\{B\}\).
For a premise other than \(B\), retain that premise and use A1;
for \(B\), use the identity derivation. For an axiom retain it
with empty support and apply A1. For MP use A2 with the two earlier
transformed lines; their support union lies in the required difference.

If a line's support does not contain \(B\), its traced subderivation
can instead be copied without any \(B\)-premise: recursively retain
only the earlier lines actually used to obtain it. Every retained
premise is in that support, while rule conditions depend only on those
unchanged support annotations. Apply A1 to the result. This observation
handles either quantified rule when \(B\) is absent from its support.

For Gen from \((C,E)\) when \(B\in E\), the original rule
condition gives \(x\notin\FV(B)\) and excludes \(x\) from
every member of \(E\setminus\{B\}\). The transformed
\(B\to C\) therefore permits Gen with its transformed support,
giving \(\forall x(B\to C)\). The registered distribution schema
and MP give \(B\to\forall xC\) with that same support. This
argument would fail if the discharged premise had a free \(x\);
the original annotated rule condition is exactly what excludes it.

For existential elimination from \((C\to G,E)\) when \(B\in E\),
the induction result is \(B\to(C\to G)\). Use the empty-support
permutation derivation of \cref{thm:p02-w02-p02-025-permutation} to obtain \(C\to(B\to G)\).
The original condition excludes \(x\) from \(B,G\) and every
member of \(E\setminus\{B\}\). The registered existential rule
therefore gives \(\exists xC\to(B\to G)\); its conclusion
\(B\to G\) has no free \(x\), and its premise support has
none. The same permutation derivation, with the two antecedents interchanged, gives
\(B\to(\exists xC\to G)\), with unchanged support. This verifies
both quantified cases and hence the strengthened induction.
Conversely append a \(B\)-premise and MP; its support is the former
support union \(\{B\}\), as required.
\end{proof}
\end{annexitem}

\begin{annexitem}{annex:p2.hfa.alpha}{thm:p02-hfa-alpha-encoding}{Literal bound renaming, closed encoding, and substitution covariance}
\begin{proof}
We first record the syntax facts used here, because the substitution
algorithm can rename a binder even when the replaced variable has no
free occurrence in its matrix. View a finite formula as its constructor
tree with each variable occurrence tagged as free with its name, or
bound by the indicated quantifier node on its path. Consistent fresh
binder renaming leaves this data unchanged. Conversely two formulas
with the same constructors, term trees, free names and binding links
can be made literally identical: visit corresponding quantifier nodes
from the root downward and rename each pair to the same name fresh
for both remaining formulas and all earlier chosen names. Each step
is an allowed fresh alpha generator. This finite construction proves
their alpha-equivalence.

The total substitution clauses preserve this binding data except for
the specified free occurrences, into which the term is inserted. When
the binder equals the replaced variable, substitution stops; when it
is different and absent from the term, the recursion crosses it; when
it occurs in the term, the explicit fresh renaming retains every old
binding link and leaves inserted term variables free. Thus inspection
of the three actual clauses gives the binding-data description, not
a capture-permitting textual convention. In particular, for \(y\)
fresh for \(A\),
\((A[y/x])[x/y]=_{\alpha}A\). It also gives the alpha-equivalence
of \(QxA\) and \(Qy(A[y/x])\) whenever \(y\) is not free in
\(QxA\), with total capture avoidance, even if that \(y\) was a
bound name inside \(A\). The common-fresh-node construction just
described reduces this assertion to the source's fresh generators.

We now derive alpha-conversion in Hilbert, without semantic completeness.
Induct on formula complexity to derive both directions for every alpha
equivalence. Reflexivity is the A1/A2 identity theorem; symmetry
interchanges the two directions, and transitivity composes them using
the propositional kernel. Negation and implication congruence use the
corresponding truth-function congruence kernels. For universal congruence,
given an empty-support theorem \(C\to G\), assume \(\forall xC\),
instantiate with \(x\) to get \(C\), use MP to get \(G\), and
apply Gen. Its only support is \(\{\forall xC\}\), which has no
free \(x\). Discharge gives \(\forall xC\to\forall xG\).
For existential congruence, compose \(C\to G\) with the registered
\(G\to\exists xG\) instance, then use existential elimination
on the resulting empty-support theorem \(C\to\exists xG\).
The variable \(x\) is free in neither the consequent nor its empty
support. This gives \(\exists xC\to\exists xG\). Repeat with
the reverse implication for the other congruence direction.

For a root universal fresh-renaming generator, put \(C'=C[y/x]\)
with \(y\notin\operatorname{Var}(C)\); the identity-name case is
trivial. From \(\forall xC\), instantiate with \(y\) to get
\(C'\), then Gen on \(y\), which is not free in the supporting
major premise. Discharge yields \(\forall xC\to\forall yC'\).
For the reverse direction, instantiate \(\forall yC'\) with \(x\).
The resulting \(C'[x/y]\) is alpha-equivalent to \(C\) by the
syntax fact above. Both matrices have lower complexity than the
quantified formula, so the induction hypothesis supplies the
empty-support implication to \(C\). MP and Gen on \(x\) now give
\(\forall xC\); \(x\) is not free in the supporting
\(\forall yC'\). Discharge gives the reverse implication.

For a root existential generator, the registered introduction schema
with term \(x\) gives \(C'[x/y]\to\exists yC'\). Compose it
with the lower-complexity theorem \(C\to C'[x/y]\), then use
existential elimination on \(x\). Its empty support and consequent
\(\exists yC'\) have no free \(x\), giving
\(\exists xC\to\exists yC'\). Conversely the introduction
instance \(C[y/x]\to\exists xC\) is literally
\(C'\to\exists xC\); existential elimination on \(y\) is
licensed by the freshness of \(y\) and empty support, and gives
the reverse direction. These generator and congruence constructions
cover the least equivalence/congruence definition of alpha-equivalence.
There is no recursive use of the claim at the quantified formula's
own complexity.

For the encoding, derive \(Q\to Q\) by A1/A2 and apply Gen to
each variable in its universal closure. Every line has empty support,
so each Gen condition holds. This constructs the closed theorem \(T\).
The formula \(F=\neg T\) is closed and is false in every classical
structure. Structural induction now proves \(\FV(h(A))=\FV(A)\):
the atom clause is unchanged, falsum contributes no free variable,
each Boolean clause takes the same union, and each quantifier removes
the same bound variable. Truth preservation follows by the same
induction, using the classical truth functions of the conjunction and
disjunction encodings and the identical quantifier domains.

Finally \(h\) preserves binding data and commutes with term insertion
at every original free occurrence. At an encoded falsum occurrence
the inserted closed subtree \(F\) has no such free occurrence;
substitution there can change only its bound names. The Boolean
encodings introduce no binders; the quantifier encodings keep their
original binding links. Applying the earlier exact substitution-clause
calculation on each side therefore gives the same constructor/term
trees, free names and binding links for \(h(A[t/x])\) and
\(h(A)[t/x]\). The common-fresh-node construction proves their
alpha-equivalence. The same argument proves preservation of an
original alpha-equivalence by \(h\). This includes unconditional
fresh-name choices made by the actual syntax algorithm.
\end{proof}
\end{annexitem}

\begin{annexitem}{annex:p2.hfa.exists-duality}{thm:p02-hfa-existential-duality}{The exact primitive-existential composition bridge}
\begin{proof}
Abbreviate \(E=\exists xA\) and \(U=\forall x\neg A\).
Universal instantiation with the term \(x\) gives
\(U\to\neg A\), since \((\neg A)[x/x]=\neg A\) literally
under the actual substitution recursion. Propositional contraposition
and double-negation introduction turn this empty-support theorem into
\(A\to\neg U\). Apply the registered existential rule on \(x\).
Its consequent \(\neg U\) has no free \(x\), and its premise
support is empty. The result is \(E\to\neg U\).

For the reverse direction, existential introduction with term \(x\)
gives the empty-support schema \(A\to E\). Contraposition gives
\(\neg E\to\neg A\). Generalize this whole formula on \(x\);
its support is empty, so Gen is licensed. The registered universal
distribution schema then gives \(\neg E\to\forall x\neg A\),
because \(x\notin\FV(\neg E)\). Contraposition gives
\(\neg U\to\neg\neg E\); compose with the propositional
double-negation-elimination kernel to get \(\neg U\to E\).
Every line has empty support. The two directions therefore remain
available under any premise context without a free-variable condition
on that unused context. This establishes precisely the source
composition claim without replacing the registered rules by a
defined-existential calculus or citing FOL completeness to prove itself.
\end{proof}
\end{annexitem}

\begin{annexitem}{annex:p2.hfa.bridges}{thm:p02-hfa-proof-bridges}{All literal Hilbert and equality-free ND rule cases}
\begin{proof}
First translate each Hilbert line to ND with open assumptions contained
in that line's annotated support. A premise is an assumption leaf.
A1--A3 have closed propositional ND derivations by the independently
proved propositional ND completeness theorem, Theorem~\ref{thm:prop-completeness}.
Substitute the relevant first-order formulas in their finite propositional
trees and freshly rename copied discharge labels; no quantified inference
is involved in that substitution. MP becomes implication elimination,
combining separately freshened copies of its two earlier trees.

The universal-instantiation schema is obtained by assuming its universal
antecedent, applying ND universal elimination with the actual total
substitution, and nonvacuously discharging that antecedent. Existential
introduction's schema is analogous: assume its displayed instance,
apply ND existential introduction and discharge. For distribution,
assume \(\forall x(C\to G)\) and then \(C\); instantiate the
universal with \(x\) and apply implication elimination to obtain
\(G\). ND universal introduction on \(x\) is legal because
neither the universal major premise nor \(C\) has free \(x\).
Discharge \(C\) and then the universal premise. Both are actually
used, so the ND nonvacuous conditions hold.

Gen translates by ND universal introduction on its earlier tree;
the Hilbert condition excludes its variable from all actual open
assumptions, since they lie in that line's support. For the Hilbert
existential rule from \((C\to G,D)\), generalize its translated
ND tree on \(x\), permitted by \(x\notin\FV(D)\).
Choose \(y\) fresh for this finite tree, \(D,C,G\), and the new
outer assumptions. Under the assumption \(\exists xC\), use
the minor assumption \(C[y/x]\). Instantiate the generalized
\(\forall x(C\to G)\) with \(y\); the result is
\(C[y/x]\to G\), because \(x\notin\FV(G)\) and \(y\)
is absent from all of \(G\), so here \(G[y/x]=G\) literally.
Application gives \(G\). ND existential elimination nonvacuously
discharges the minor assumption, with \(y\) absent from every
surviving assumption, the major existential and the conclusion.
Finally discharge the existential assumption to obtain
\(\exists xC\to G\). This checks the stricter ND witness absence,
not merely the Hilbert free-variable restriction. These are all
Hilbert rows and introduce no equality rule.

For the reverse direction, fix the closed \(T\), write \(F=\neg T\),
and induct over the whole equality-free ND tree. Each translated line
will have support contained in the images of its permitted open
assumption formulas. Assumptions are translated premise lines.
The following propositional constructions use only the kernel lemma,
MP and the tracked discharge lemma, so their support accounting is
by union and removal, not an unlicensed new inference rule.

For conjunction introduction, from \(h(C),h(G)\) use the tautology
\(h(C)\to(h(G)\to\neg(h(C)\to\neg h(G)))\).
Each conjunction elimination uses the corresponding projection
tautology from \(\neg(h(C)\to\neg h(G))\).
Each disjunction introduction uses its injection tautology into
\(\neg h(C)\to h(G)\). For disjunction elimination, discharge
the two translated minor assumptions to obtain \(h(C)\to h(H)\)
and \(h(G)\to h(H)\). The classical Cases tautology combines
these with the translated major \(\neg h(C)\to h(G)\), giving
\(h(H)\); its support is contained in the three surviving contexts.
Implication introduction uses tracked premise discharge, and implication
elimination is literal MP.

For negation introduction, the minor translation gives \(F\) under
\(h(C)\); discharge gives \(h(C)\to\neg T\). Combine this
with the empty-support theorem \(T\) using the tautology
\(T\to((h(C)\to\neg T)\to\neg h(C))\).
Negation elimination combines \(h(C),\neg h(C)\) with the
contradiction-to-\(F\) tautology. Falsity elimination combines
\(F=\neg T\) and the theorem \(T\) with explosion's classical
truth-function kernel to obtain the requested translated formula.
For RAA, the minor translated contradiction under \(\neg h(C)\)
gives \(\neg h(C)\to\neg T\) by discharge. Use
\(T\to((\neg h(C)\to\neg T)\to h(C))\) and MP.
This proves the classical case without pretending that \(F\) is a
primitive Hilbert falsity or importing a new RAA rule. The ND input
may require every discharge to be nonempty; those actual tree classes
are present. Hilbert's deduction construction needs no added unused
ND assumption or repetition rule.

For universal elimination, use the registered schema to obtain
\(h(C)[t/x]\) from \(\forall xh(C)\). The alpha/substitution
lemma supplies the empty-support implication to \(h(C[t/x])\),
so MP restores the exact requested conclusion string. Existential
introduction uses the reverse implication between those two instances,
then the registered existential-introduction schema. Universal
introduction applies Gen directly to \(h(C)\): equality of free
variable sets ensures that its variable is free in none of the
translated surviving premises.

For ND existential elimination, let the witness be \(y\), the major
be \(\exists xC\), and the conclusion be \(H\).
The minor induction result is from the translated other assumptions
and \(h(C[y/x])\). Discharge that witness formula. Substitution
covariance and the propositional implication-congruence kernel give
\(h(C)[y/x]\to h(H)\), with support in the translated surviving
minor context. Apply the registered existential rule on \(y\).
It is licensed because \(y\) is free in neither \(h(H)\) nor
that support, by free-variable preservation and the original ND
condition. This gives
\(\exists y(h(C)[y/x])\to h(H)\).
The original witness condition puts \(y\) outside the free variables
of \(\exists xC\); hence
\(\exists y(h(C)[y/x])=_{\alpha}\exists xh(C)\), including
any bound \(y\)'s introduced inside the closed encoding. The alpha
lemma supplies the appropriate empty-support implication. Compose
and apply the translated major premise to get exactly \(h(H)\).
Thus the original major/minor support union is preserved. There is
no leaked witness premise and no illicit Gen on its open minor
assumption.

Finally an ND alpha step translates by alpha-preservation of \(h\)
and the derived Hilbert alpha implication. These exhaust all
propositional, quantifier and alpha rows of the equality-free restriction.
The ND equality introduction/elimination rows are excluded throughout
the input tree, not merely because its end formula lacks equality.
Finite concatenation of the constructed Hilbert sequences completes
the second translation.
\end{proof}
\end{annexitem}

\begin{annexitem}{annex:p2.hfa.adequacy}{thm:p02-hfa-adequacy}{Finite support, exact-language adequacy, and the signature boundaries}
\begin{proof}
Soundness was proved directly for the literal sheet. For the converse
with arbitrary \(\Gamma\), use the inherited strong Henkin completeness
theorem in its equality-free specialization. The actual specialization
is identified in Annex~\ref{annex:p7.calculi.4}: it omits the equality
rules, congruence step and term quotient, using the literal term model;
its connective, quantifier, witness and truth arguments remain those
of the proved construction. It therefore supplies an equality-free
ND tree, not an equality-bearing tree whose endpoint happens to omit
equality. The inherited arbitrary-signature construction assumes Choice,
so this step is in classical ZFC as stated.

Choose an equality-free atom occurring in the fragment conclusion
\(A\), and construct the closed \(T\) from it. Every fragment
formula contains such an atom, because its bottom-free grammar starts
with atoms. Translate the complete ND tree by the second proof bridge.
The map \(h\) is literally the identity on all fragment premises and
on \(A\); it changes only extra full-language constructors occurring
inside the ND proof. Thus the result is a derivation of exactly \(A\)
from exactly the allowed set \(\Gamma\) on the registered sheet.
Its finite sequence has finite support, without assuming that the
entire premise set is finite. This proves strong adequacy noncircularly.

For finite \(\Gamma\), let \(\mathcal L_0\) contain the finitely
many nonlogical symbols in its formulas and \(A\). A structure for
\(\mathcal L_0\) expands to the ambient signature: choose one
element of its nonempty domain, interpret every unused function and
constant uniformly by that element, and every unused relation by the
empty relation. This is one uniform construction, not an infinite
family of choices. Satisfaction of input formulas is unchanged by
reduct invariance. Therefore validity in the ambient signature is
validity in this finite signature. The supplied countable-language
Henkin route, Corollary~\ref{cor:fol-completeness-countable}, with the
same equality-free specialization gives the ND tree in classical ZF.
The proof bridges use only finite syntax and fresh variables and add
no Choice assumption. This gives the finite-input ZF statement even
for arbitrary ambient signatures. For arbitrary \(\Gamma\) in
a supplied countable signature, use that same countable-language
strong completeness directly. No effective premise membership or
decision procedure for unrestricted FOL is inferred.

For the full-language equivalence, truth preservation gives
\(\Gamma\entails A\) iff \(h[\Gamma]\entails h(A)\) in the
same signature whenever the encoding's atom exists. Apply the
just-proved fragment adequacy, or equivalently translate the complete
ND proof directly. In the atomless case adjoin the expressly stated
nullary predicate. Every original structure has an expansion (give
that predicate any truth value), every expanded structure has an
original reduct, and \(T\) is a closed tautology in every expansion.
Truth preservation is therefore independent of the added predicate's
value; reduct/expansion give both directions of the equivalence.
This proves the exceptional full-language translation without claiming
that an atomless literal fragment suddenly contains falsity.

If logical equality were allowed as an additional atom but assigned
identity semantics without adding identity rules, the claim would
be a different theorem. Indeed the same rule-soundness proof holds
when that symbol is interpreted as an arbitrary binary relation;
interpreting it empty in a nonempty structure falsifies \(x=x\).
Thus a reflexivity conclusion would not be derivable from the
unchanged rules in that altered language. On the actual sheet it
is outside the equality-free grammar, so this observation is an
extension boundary, not an alleged counterexample to its proved
adequacy. The equality-bearing ND/LK proofs are not imported.
\end{proof}
\end{annexitem}
